\chapter{Fever in Children}

\section*{Definition and Overview}
Fever is a rise in body temperature above the normal range, generally regarded as a temperature of 38 degrees Celsius or higher, and is one of the most common reasons a child is brought to a doctor. Fever is not an illness in itself but a sign that the immune system is fighting an infection. Most fevers in children are caused by minor viral infections and resolve within a few days, but fever can occasionally be the first sign of a serious bacterial infection such as pneumonia, a urinary tract infection, or meningitis. The challenge for parents and clinicians is to distinguish the common, harmless fever from the rare but dangerous one. How the child appears and behaves matters far more than the height of the temperature.

\section*{Epidemiology}
Fever is extremely common in childhood. Most children will have several fevers before they are five years old, usually from viral respiratory infections, ear infections, and gastroenteritis. Fevers are most frequent in infancy and early childhood, partly because the immune system is still developing and partly because young children are exposed to many viruses. Serious bacterial infections are uncommon but are more likely in very young babies, particularly those under three months, and in children with certain underlying conditions. Since the introduction of conjugate vaccines, diseases such as pneumococcal and meningococcal infection, which can present with fever, have become much rarer.

\section*{Aetiology and Pathophysiology}
\begin{itemize}
    \item \textbf{Viral infections:} by far the most common cause, including colds, influenza, croup, bronchiolitis, and childhood exanthems such as roseola
    \item \textbf{Bacterial infections:} include ear infections, tonsillitis, urinary tract infection, pneumonia, and, rarely, meningitis or sepsis
    \item \textbf{Post-vaccination fever:} a normal immune response to immunisation that settles within a day or two
    \item \textbf{Non-infectious causes:} teething can cause mild temperature elevation, and overheating from excessive clothing or a hot environment can raise temperature
    \item \textbf{Mechanism:} infection triggers the release of pyrogens that act on the hypothalamus, resetting the body's thermostat upwards to make the environment less favourable for microbes and to enhance immune responses
    \item \textbf{Fever itself is not dangerous:} the height of the temperature alone does not predict severity, and fever does not cause brain damage; what matters is the underlying illness
\end{itemize}

\section*{Clinical Features}
\begin{itemize}
    \item A measured temperature of 38 degrees Celsius or higher
    \item The child may feel hot to the touch and appear flushed
    \item Reduced appetite, fussiness, or being less active than usual
    \item Accompanying symptoms of the underlying illness, such as cough, runny nose, sore ears, vomiting, or diarrhoea
    \item \textbf{Warning features of serious illness:} lethargy or difficulty rousing, floppiness, rapid or laboured breathing, a rash that does not fade when pressed, a bulging fontanelle in a baby, and fits
    \item \textbf{Febrile seizure:} a fit triggered by fever, most common between 6 months and 5 years of age
\end{itemize}

\section*{Differential Diagnosis}
\begin{itemize}
    \item Viral upper respiratory tract infection and influenza-like illness
    \item Acute otitis media and tonsillitis
    \item Gastroenteritis with dehydration
    \item Urinary tract infection and pyelonephritis
    \item Pneumonia and bronchiolitis
    \item Meningitis and septicaemia -- rare but must always be considered
    \item Roseola infantum and other childhood viral rashes
    \item Post-vaccination fever
    \item Heat exhaustion and overheating
\end{itemize}

\section*{When to Seek Medical Care}
See a doctor if a child has a fever that persists beyond about three days, if the child seems unusually unwell, if fluids are not being taken well, or if there are concerns about the child's breathing or behaviour.

\medskip
\textbf{Call 000 (or local emergency services) for:}
\begin{itemize}
    \item Fever in a baby under three months of age
    \item A child who is drowsy, floppy, difficult to rouse, or unusually irritable
    \item A rash that does not fade when a glass is pressed against it
    \item Difficulty breathing, fast breathing, or signs of chest infection
    \item A febrile seizure or fit
    \item A child who is very pale, mottled, cold, or appears seriously unwell
    \item An inability to keep down fluids, with signs of dehydration such as reduced urine output
\end{itemize}

\section*{Diagnosis and Investigations}
\begin{itemize}
    \item \textbf{History:} the height and duration of the fever, how it was measured, other symptoms, vaccination status, fluid intake, and urine output
    \item \textbf{Examination:} temperature, breathing rate and effort, heart rate, hydration, skin (including checking for a non-blanching rash), ears, throat, and overall alertness and interaction
    \item \textbf{Assessment of risk:} the age of the child and the presence of warning features determine whether tests are needed
    \item \textbf{Urine testing:} a dipstick and culture are commonly performed in young children to exclude a urinary tract infection
    \item \textbf{Blood tests:} full blood count, inflammatory markers, and blood cultures in children who look unwell or are at higher risk of serious bacterial infection
    \item \textbf{Chest X-ray:} if pneumonia is suspected
    \item \textbf{Lumbar puncture:} performed when meningitis is suspected, to examine the fluid around the brain and spinal cord
    \item \textbf{Observation:} some children are observed in hospital to see how their condition evolves
\end{itemize}

\section*{Management}
\begin{itemize}
    \item \textbf{Reassurance and monitoring:} most fevers settle in two to three days; parents should keep the child comfortable and watch for warning signs
    \item \textbf{Antipyretics:} paracetamol or ibuprofen, given only if the child is distressed by the fever, in a weight-appropriate dose and never both together unless advised
    \item \textbf{Fluids:} encourage regular small drinks to prevent dehydration; oral rehydration solution helps after vomiting or diarrhoea
    \item \textbf{Dress lightly:} avoid over-wrapping, and do not sponge with cold water, which is unnecessary and uncomfortable
    \item \textbf{Treat the cause:} antibiotics for bacterial infections such as ear, urine, or chest infections, prescribed by a doctor
    \item \textbf{Hospital care:} for young babies, seriously unwell children, or those needing intravenous fluids or antibiotics
    \item \textbf{Do not give aspirin} to children or teenagers under 16 years, because of the risk of Reye's syndrome
\end{itemize}

\section*{Complications}
\begin{itemize}
    \item Dehydration from reduced intake and increased fluid loss
    \item Febrile seizures, which look alarming but in most cases cause no lasting harm
    \item Progression of an unrecognised serious bacterial infection, such as meningitis or sepsis
    \item Worsening of the underlying illness, such as pneumonia, if care is delayed
    \item Anxiety and lost sleep for both child and parents
\end{itemize}

\section*{Prognosis and Outcomes}
The great majority of childhood fevers are caused by self-limiting viral infections and settle within a few days with supportive care. Serious bacterial infections are uncommon and, when recognised and treated promptly, usually respond well to antibiotics. Febrile seizures generally resolve without long-term consequences. The outlook is excellent for children who look well between fevers, drink and pass urine normally, and have no warning features. Parents should trust their instincts: if a child seems seriously unwell, they should seek urgent medical advice.

\section*{Prevention and Self-Care}
\begin{itemize}
    \item Keep childhood vaccinations up to date, including against pneumococcus, meningococcus, and influenza
    \item Practise regular hand washing to reduce the spread of infections
    \item Measure temperature accurately and record how the child behaves, not just the number
    \item Keep a child with fever comfortable, lightly dressed, and well hydrated
    \item Use paracetamol or ibuprofen only when the child is distressed, in the correct weight-based dose
    \item Do not wake a sleeping child to give fever medicine
    \item Call a doctor for fever in a baby under three months without waiting
    \item Seek help early if the child is not drinking, has reduced urine output, or seems to be getting worse
\end{itemize}

\section*{Sources and Further Reading}
The following reputable sources provide further information on fever in children:
\begin{itemize}
    \item NHS. \textit{Health A--Z: fever in children information.} https://www.nhs.uk/conditions/
    \item NICE. \textit{Clinical guideline on fever in under 5s.} https://www.nice.org.uk/
    \item Mayo Clinic. \textit{Symptoms and causes of fever in children.} https://www.mayoclinic.org/
    \item Centers for Disease Control and Prevention. \textit{Childhood fever and illness resources.} https://www.cdc.gov/
    \item MedlinePlus. \textit{Health topics on children's health.} https://medlineplus.gov/
    \item World Health Organization. \textit{Child health and fever management resources.} https://www.who.int/
\end{itemize}
